% 张三慧大学物理学第四版第3章3.1--3.20独立解答
% 依据教材力学热学PDF116--117页；配套辅导书第3章同题号比对待继续。
% 原书图片的条件在 question.tex 记录；未经外审或版权许可核验。
\section*{第3章\quad 动量与角动量：独立解析}
\subsection*{章末习题3.1--3.20}

\paragraph{习题 3.1}
小球所受弹力的冲量等于该力在时间内的积分，
\[
 \boldsymbol J=\int_0^{\pi/(2\omega)}
 (-kA\cos\omega t)\,dt\,\boldsymbol i
 =-\frac{kA}{\omega}
 \left[\sin\omega t\right]_0^{\pi/(2\omega)}\boldsymbol i
 =\boxed{-\frac{kA}{\omega}\boldsymbol i}.
\]
负号表示向负 \(x\) 方向，量纲为
\(\mathrm{N\,s}\)。结果还须符合动量定理
\(\boldsymbol J=m[\boldsymbol v(t_2)-\boldsymbol v(t_1)]\)
（若 \(k=m\omega^2\)，两者相同）。

\paragraph{习题 3.2}
匀速圆周运动的速率固定为 \(v\)，经历 \(1/4\)
周期时速度方向转过 \(90^\circ\)。
由动量定理，
\[
 |\boldsymbol J|=m|\boldsymbol v_f-\boldsymbol v_i|
 =m\sqrt{v^2+v^2}
 =\sqrt2(0.050)(20)
 \approx\boxed{1.41\,\mathrm{N\,s}}.
\]
冲量的方向沿两动量矢量之差而非简单径向固定方向；
向心力方向随时间变化，不能用初始向心力乘时间。

\paragraph{习题 3.3}
从最高点自由落下 \(h=5.0\,\mathrm m\)，
忽略空气阻力，触水瞬间向下速率
\(v=\sqrt{2gh}\approx9.90\,\mathrm{m/s}\)。
以向上为正，触水后
\(\Delta t=0.30\,\mathrm s\) 内速度由 \(-v\) 变为近零。
若水对人的平均作用力为 \(F_{\rm water}\)，
\[
 (F_{\rm water}-mg)\Delta t
 =m(0-(-v)).
\]
解得
\[
 \boxed{F_{\rm water}=m\left(g+\frac{\sqrt{2gh}}{\Delta t}\right)
 \approx150(9.80+9.8995/0.30)
 \approx6.42\times10^3\,\mathrm N}.
\]
该值是平均支持/阻力而非水击瞬时峰值，
忽略人体部分浸入时的浮力单独分解。

\paragraph{习题 3.4}
子弹发射率
\(\dot N=120/60=2\,\mathrm{s^{-1}}\)，
单发质量 \(m=7.90\times10^{-3}\,\mathrm{kg}\)。
忽略枪口气体动量，
枪托对肩部的平均反冲力大小
\[
 \boxed{\bar F=\dot Nmv
 =2(0.00790)(735)\approx11.6\,\mathrm N}.
\]
方向为枪射击方向的反向。
原题提到“以分钟计”时，应将每分钟发射数换算成
每秒数，否则会有60倍错误。

\paragraph{习题 3.5}
空气从风速 \(v=38\,\mathrm{m/s}\) 降为零，
垂直车厢侧面的迎风面积
\(A=(22.5)(2.62)=58.95\,\mathrm{m^2}\)。
在流场全截面完全停下的简化假设下，
每秒撞击空气质量为
\(\dot m=\rho Av\)，
平均动量流变化给出
\[
 \boxed{F=\dot mv=\rho Av^2
 =(1.29)(22.5)(2.62)(38)^2
 \approx1.10\times10^5\,\mathrm N}.
\]
力沿风吹向车厢的方向。
若采用 \(F=\frac12C_D\rho Av^2\) 的
空气动力阻力模型，必须另有 \(C_D\)；
不能在原题“碰后停下”的动量模型中
无依据地额外乘 \(1/2\)。

\paragraph{习题 3.6}
将弯管前的水速设为 \(v\boldsymbol i\)，
转弯后为 \(v\boldsymbol j\)，
\(\boldsymbol i,\boldsymbol j\) 互相垂直。
稳态水流动量定律（仅动量通量的净改变）为
\[
 \boldsymbol F_{\rm pipe\to water}
 =\dot m(\boldsymbol v_{\rm out}-\boldsymbol v_{\rm in})
 =\dot mv(\boldsymbol j-\boldsymbol i).
\]
因此水对管的反作用力为
\[
 \boxed{\boldsymbol F_{\rm water\to pipe}
 =\dot mv(\boldsymbol i-\boldsymbol j)},\qquad
 \boxed{|F|=\sqrt2\,\dot mv
 =\sqrt2(3\times10^3)(10)
 \approx4.24\times10^4\,\mathrm N}.
\]
方向相对于入口流向偏向出口速度的反方向
\(45^\circ\)。
\textbf{模型限制：}真实管壁总载荷还包括
进口、出口压力在各截面造成的作用；
题中未给压力值，本答案是动量流转向所需的净力贡献，
不能把它无条件认作实际全管压强载荷。

\paragraph{习题 3.7}
取已提起的链长为 \(l\)，单位长度质量
\(\rho_l\)，速度恒定 \(v_0\)。
运动链段重量为 \(\rho_l lg\)。
在时间 \(dt\) 内有质量
\(dm=\rho_lv_0dt\) 从桌面上的静止态
进入竖直运动的链段，获得动量
\(v_0dm\)，故单位时间的动量输入为
\[
 \frac{d p_{\rm entrained}}{dt}
 =v_0\frac{dm}{dt}=\rho_lv_0^2.
\]
由随质量增长的系统动量平衡
\(F-\rho_l lg=\rho_lv_0^2\)，得
\[
 \boxed{F=\rho_l lg+\rho_lv_0^2}.
\]
当 \(v_0\to0\) 时恢复准静态吊起质量
\(\rho_l l\) 的重力。

\paragraph{习题 3.8}
设链的落下长度为 \(x\)，
单位长度质量 \(\rho_l\)，从静止松开，
其着地瞬间速度由
\(v^2=2gx\) 给出。
桌上已经静止堆积的部分重量为
\(\rho_l gx\)，此外
连续到达的质量流率
\(\dot m=\rho_lv\) 由速度 \(v\)
变为零，产生附加向下作用力
\(\dot mv=\rho_lv^2=2\rho_lgx\)。
故桌面的总受压为
\[
 \boxed{N=\rho_lgx+\rho_lv^2
 =3\rho_lgx=3W_{\rm deposited}}.
\]
该结论适用于理想完全柔软链逐渐落地，
不包括最后端点着地时的瞬态冲击。

\paragraph{习题 3.9}
原核初始静止，初动量为零。
令电子动量 \(\boldsymbol p_e=p_1\boldsymbol i\)，
中微子动量 \(\boldsymbol p_\nu=p_2\boldsymbol j\)，
两者正交。
则由动量守恒，反冲原子核动量为
\[
 \boldsymbol p_{\rm recoil}=-
 (\boldsymbol p_e+\boldsymbol p_\nu).
\]
大小为
\[
 \boxed{p_{\rm recoil}
 =\sqrt{p_1^2+p_2^2}
 =\sqrt{(9.22)^2+(5.33)^2}\times10^{-21}
 \approx1.065\times10^{-20}\,\mathrm{kg\,m/s}}.
\]
方向与电子动量方向反向偏向中微子反方向
\(\beta=\arctan(p_2/p_1)\approx30.0^\circ\)。
这里动量守恒不要求已知原子核的质量。

\paragraph{习题 3.10}
以火箭飞行方向为正。
两部分分离瞬间相互作用内力的冲量相反，
总动量不变：
\[
 m_1v_1+m_2v_2=(m_1+m_2)v_0.
\]
假定仪器舱相对火箭壳向前分离，
\(v_2-v_1=u=910\,\mathrm{m/s}\)。
联立得
\[
 \boxed{v_1=v_0-\frac{m_2}{m_1+m_2}u
 \approx7289.8\,\mathrm{m/s}},
\quad
 \boxed{v_2=v_0+\frac{m_1}{m_1+m_2}u
 \approx8199.8\,\mathrm{m/s}}.
\]
检查：两速度差恰为 \(910\,\mathrm{m/s}\)，
且质量加权平均速度为 \(7600\,\mathrm{m/s}\)。
如果分离方向与飞行方向相反，
相应的速度符号需按真实几何调整。

\paragraph{习题 3.11}
设相撞后两车粘在一起（质量 \(2m\)），
沿地面因摩擦减速 \(\mu_kg\)，
滑行 \(s=25\,\mathrm m\) 后停止。
由功和能，
\[
 V=\sqrt{2\mu_kgs}
 =\sqrt{2(0.8)(9.8)(25)}
 \approx19.80\,\mathrm{m/s}.
\]
相撞前两车速度互相垂直，
记其速率为 \(u_1,u_2\)。
动量守恒的矢量平方给出
\[
 (m u_1)^2+(m u_2)^2=(2mV)^2
 \quad\Rightarrow\quad
 \boxed{u_1^2+u_2^2=4V^2\approx1568}.
\]
若两车速度均不超过 \(14\,\mathrm{m/s}\)，
左端至多 \(2(14)^2=392\)，显然矛盾。
因此两个司机\emph{不可能同时}守速限；
至少有一方超速。不能由此单独断言两人各自都超速。

\paragraph{习题 3.12}
设废气喷出前总质量 \(M=6090\,\mathrm{kg}\)，
向前速度 \(v_0=105\,\mathrm{m/s}\)；
喷出废气质量 \(\Delta m=80.0\,\mathrm{kg}\)，
其相对于探测器的速度为 \(-u\)、
\(u=253\,\mathrm{m/s}\)。
喷出后探测器速度为 \(v\)，
废气在同一惯性系中速度为 \(v-u\)。
动量守恒
\[
 Mv_0=(M-\Delta m)v+\Delta m(v-u)
 =Mv-\Delta m u.
\]
故
\[
 \boxed{v=v_0+\frac{\Delta m}{M}u
 =105+\frac{80.0\times253}{6090}
 \approx108.32\,\mathrm{m/s}}.
\]
若质量值或喷气相对速度的参考系不同，
应重新从原题核验，不应沿用本式的分母。

\paragraph{习题 3.13}
设火箭无外力，初始静止，初质量 \(m_0\)，
喷气相对于火箭速度大小恒为 \(u\)。
齐奥尔科夫斯基火箭方程为
\[
 v=u\ln\frac{m_0}{m}.
\]
后喷气体在外部惯性系中静止需
\(v-u=0\)，即
\[
 \ln\frac{m_0}{m}=1,\qquad
 \frac{m}{m_0}=e^{-1}.
\]
故燃烧掉的质量占初质量的比例
\[
 \boxed{\frac{m_0-m}{m_0}=1-e^{-1}
 \approx0.6321}.
\]
此处“喷气静止”是指相对初始静止惯性系的瞬时
废气速度为零，并非所有已喷出的气体同时静止。

\paragraph{习题 3.14}
设初质量 \(m_0=2.72\times10^6\,\mathrm{kg}\)，
质量燃烧率
\(\dot m=1.29\times10^3\,\mathrm{kg/s}\)，
喷气相对速率 \(u\)，常量重力
\(g=9.80\,\mathrm{m/s^2}\)。
瞬时推力 \(T=\dot m\,u\)。
(1) 刚起飞临界条件 \(T=m_0g\)，故
\[
 \boxed{u_{\min}=\frac{m_0g}{\dot m}
 \approx2.07\times10^4\,\mathrm{m/s}}.
\]
(2) 在竖直上升且忽略阻力时
\(m(t)=m_0-\dot mt\)，
\[
 \frac{dv}{dt}=\frac{u\dot m}{m(t)}-g.
\]
以 \(v(0)=0\) 积分得
\[
 v(t)=u\ln\frac{m_0}{m_0-\dot mt}-gt.
\]
所给燃烧时间 \(t_b=155\,\mathrm s\) 后，
剩余质量
\(m_b=2.72\times10^6-(1.29\times10^3)(155)
=2.52005\times10^6\,\mathrm{kg}\)。
因 \(u\dot m/m(t)>g\) 全程成立，
速度最大值在刚燃尽时：
\[
 \boxed{v_{\max}
 =(5.50\times10^4)\ln\frac{2.72\times10^6}{2.52005\times10^6}
 -(9.80)(155)
 \approx2.68\times10^3\,\mathrm{m/s}}.
\]
(3) 继续对 \(v(t)\) 积分：
\[
 h(t)=u\left[t-\frac{m(t)}{\dot m}
 \ln\frac{m_0}{m(t)}\right]
 -\frac12gt^2.
\]
因此
\[
 \boxed{h(t_b)\approx2.04\times10^5\,\mathrm m}.
\]
\textbf{与配套辅导书的实质差异：}
配套《学习辅导与习题解答》PDF物理页53（书上印刷页44）
的第3.14题在写出积分表达式后给出 \(h=172\,\mathrm{km}\)。
使用其同页所示的 \(m_0,\dot m,u,T,g\) 直接积分
\[
 h=\int_0^{155}\left[
 (5.50\times10^4)\ln
 \frac{2.72\times10^6}{2.72\times10^6-1.29\times10^3t}
 -9.8t\right]dt
 \approx203.59\,\mathrm{km},
\]
并可用独立数值积分交叉复核。
\textbf{172 km与203.59 km不相符；此处保留原辅导书结论并
将差异标记为待审，不得静默替换。}
\textbf{模型限定：}在这个高度重力不可能严格恒定，
空气阻力也不能一直忽略；
数值只是照原题理想化近似演算，不是实际火箭工程轨迹。

\paragraph{习题 3.15}
令飞机向前为正，进气空气相对于地面近似静止，
其质量流率 \(\dot m_a=75\,\mathrm{kg/s}\)，
燃料在飞机内初始随飞机以 \(v=210\,\mathrm{m/s}\) 运动，
燃料质量流率 \(\dot m_f=3.0\,\mathrm{kg/s}\)。
尾气以相对于飞机 \(-u=-490\,\mathrm{m/s}\)
喷出，地面系的尾气速度是 \(v-u=-280\,\mathrm{m/s}\)。
发动机向后输出的流体总动量变化，
由反作用得飞机推力
\[
 F=\dot m_a(0)+\dot m_fv
 -(\dot m_a+\dot m_f)(v-u)
 =(\dot m_a+\dot m_f)u-\dot m_av.
\]
即
\[
 \boxed{F=78(490)-75(210)
 =2.247\times10^4\,\mathrm N}.
\]
若把进气初动量误当成随飞机速度 \(v\)，
会漏掉“冲压进气”造成的动量差。

\paragraph{习题 3.16}
两根 O--H 键等长
\(l=0.0958\,\mathrm{nm}\)，夹角 \(105^\circ\)。
水分子质心必在角平分线上。
以氧原子核位置为原点、向两氢原子中间为正，
氢原子到该轴的投影均为
\(l\cos(105^\circ/2)\)。
取相对原子质量 \(m_H:m_O\simeq1:16\)，有
\[
 x_{\rm CM}=
 \frac{2m_Hl\cos52.5^\circ}{m_O+2m_H}
 =\frac{2(0.0958)\cos52.5^\circ}{18}\,\mathrm{nm}
 \approx\boxed{0.00648\,\mathrm{nm}}.
\]
因此质心在从 O 核指向两个 H 核夹角平分线
的方向上，离 O 核约 \(6.48\times10^{-3}\,\mathrm{nm}\)。

\paragraph{习题 3.17}
均匀半圆形薄板半径记为 \(R\)，
取直径沿 \(x\) 轴，材料在 \(y\ge0\)。
由于左右镜像对称，\(x_{\rm CM}=0\)。
用极坐标 \(y=r\sin\theta\)，
区域 \(0\le r\le R,\ 0\le\theta\le\pi\)，
质心纵坐标
\[
 y_{\rm CM}=
 \frac{\int_0^\pi\int_0^R
 r\sin\theta\cdot r\,dr\,d\theta}
 {\int_0^\pi\int_0^Rr\,dr\,d\theta}
 =\frac{(R^3/3)\cdot2}{(\pi R^2/2)}
 =\boxed{\frac{4R}{3\pi}}.
\]
位置沿半圆对称轴朝向弧面。
\textbf{注意：}这里是“半圆形薄板”，
不是“半圆细线”；半圆细线的质心为 \(2R/\pi\)。

\paragraph{习题 3.18}
取边长 \(a\) 的均匀立方体，
以底面为 \(z=0\)，顶面为 \(z=a\)。
图3.38显示一根沿前后方向贯穿立方体的
半径 \(a/4\) 圆柱形孔，孔在下半部中央，
其截面与底面相切，故孔轴高度 \(z_h=a/4\)。
原立方体质量
\(M=\rho a^3\)，孔质量
\(m_h=\rho\pi(a/4)^2a=\rho\pi a^3/16\)。
利用“完整体减去挖空部分”的质心公式，
\[
 z_{\rm remain}
 =\frac{M(a/2)-m_h(a/4)}{M-m_h}
 =a\frac{1/2-\pi/64}{1-\pi/16}
 \approx\boxed{0.5611a}.
\]
配套辅导书PDF物理页54把质心相对\emph{立方体几何中心}
表示为向上 \(0.061a\)，与这里从\emph{底面}
量起的 \(0.5611a\) 一致，二者不是解答冲突。
其平面位置仍在左右及前后中心线的交点。
\textbf{复核依赖：}若图3.38实际孔轴高度
不是 \(a/4\)（即孔不与底面相切），
应把公式中的 \(a/4\) 换为真实孔轴高度 \(z_h\)；
本数值需图示最终审定。

\paragraph{习题 3.19}
两石子从同一处相继释放，取向下为正，
以第一石子释放时为 \(t=0\)，
两质量分别
\(m_1=0.020\,\mathrm{kg}\)、
\(m_2=0.050\,\mathrm{kg}\)。
当 \(t>1\,\mathrm s\) 且均未落地时，
\[
 v_1=gt,\qquad v_2=g(t-1).
\]
质心速度为
\[
 v_C=\frac{m_1v_1+m_2v_2}{m_1+m_2}
 =g\left[t-\frac{m_2}{m_1+m_2}(1\,\mathrm s)\right]
 =\boxed{9.80(t-5/7)\,\mathrm{m/s}},
\]
式中 \(t\) 以秒为数值单位。
两石子都仅受重力，因此
\[
 \boxed{a_C=\frac{m_1g+m_2g}{m_1+m_2}
 =9.80\,\mathrm{m/s^2}}
\]
方向向下。若任一石子已落地，
动量系统的外力条件变化，不能继续套用。

\paragraph{习题 3.20}
太阳引力为中心力，
对太阳参考点的力矩为零，
所以彗星单位质量的角动量
\(h=|\boldsymbol r\times\boldsymbol v|\)
在各时刻守恒。
近日点与远日点的径向速度分量均为零，
即速度与径矢垂直：
\[
 r_1v_1=r_2v_2.
\]
故
\[
 \boxed{r_2=\frac{r_1v_1}{v_2}
 =\frac{(8.75\times10^{10})(5.46\times10^4)}
 {9.08\times10^2}
 \approx5.26\times10^{12}\,\mathrm m}.
\]
检查：\(r_2>r_1\)，且远日点速率更低，
符合中心力角动量守恒。


% 第3章末尾的4道习题：PDF物理页118。
\subsection*{续录习题3.21--3.24：独立解析}

\paragraph{习题 3.21}
采用题设圆轨道近似，
取月球质量 \(m_{\rm M}\simeq7.35\times10^{22}\,\mathrm{kg}\)，
地月中心距 \(r\simeq3.844\times10^8\,\mathrm m\)；
周期 \(T=27.3(86400)\,\mathrm s\)。
月球速度沿切线，故对地心的角动量大小
\[
 \boxed{L=m_{\rm M}rv
 =\frac{2\pi m_{\rm M}r^2}{T}
 \approx2.89\times10^{34}\,\mathrm{kg\,m^2/s}}.
\]
月地连线每单位时间扫过面积为
\[
 \boxed{\frac{dA}{dt}=\frac{\pi r^2}{T}
 \approx1.97\times10^{11}\,\mathrm{m^2/s}}.
\]
检查 \(dA/dt=L/(2m_{\rm M})\)；
这一关系对任意中心力轨道成立，此处用圆轨道代数值。

\paragraph{习题 3.22}
设地球平均半径近似 \(R_\oplus=6371\,\mathrm{km}\)
（具体数值随教材数据表略有不同）。
近、远地点到地心的距离
\[
 r_p=R_\oplus+205.5\,\mathrm{km},\qquad
 r_a=R_\oplus+35835.7\,\mathrm{km}.
\]
万有引力是中心力，角动量守恒；
近远地点速度都沿切线，所以
\[
 r_pv_p=r_av_a
 \quad\Rightarrow\quad
 \boxed{v_a=\frac{r_p}{r_a}v_p
 \approx1.59\,\mathrm{km/s}}.
\]
椭圆的长半轴
\(a_e=(r_p+r_a)/2\)，短半轴
\(b_e=\sqrt{r_pr_a}\)。
按开普勒第二定律，掠面速度
\(\dot A=r_pv_p/2\) 不变，整圈面积
\(\pi a_eb_e\)。
因此
\[
 \boxed{T=\frac{\pi a_eb_e}{r_pv_p/2}
 =\frac{2\pi a_eb_e}{r_pv_p}
 \approx3.81\times10^4\,\mathrm s
 \approx10.57\,\mathrm h}.
\]
上述推导不需要已知卫星质量；
近地点与远地点高度都必须先加地球半径再取比。

\paragraph{习题 3.23}
若细绳始终沿半径向内拉，绳张力对小方块
绕桌孔中心的力矩为零；
桌面光滑也没有其他水平力矩。
故角动量守恒
\[
 m r_0v_0=mrv,\qquad
 \boxed{v=v_0\,\frac{r_0}{r}}.
\]
因 \(r\) 变小，速率变大；
增加的动能来自拉绳外力做功，
所以不能将机械动能误认为恒定。

\paragraph{习题 3.24}
选两质点连线为 \(x\) 轴，原两质点位于
\(x=-a/2,+a/2\)，各质量 \(m\)，
其质心 \(O\) 静止且绕 \(O\) 以 \(\omega\) 旋转。
第三个静止质量 \(m\) 的质点与 \(x=+a/2\) 端碰撞。
(1) 碰前瞬间三个质点总体质心为
\[
 x_C=\frac{m(-a/2)+m(a/2)+m(a/2)}{3m}
 =\boxed{\frac a6}
\]
（从原杆中心向碰撞端）。
杆两端动量矢量互相抵消，第三点动量为零，
故总动量零，\(\boxed{\boldsymbol v_C=\boldsymbol0}\)。
(2) 因总动量为零，
对新质心 \(C\) 的角动量与对原质心 \(O\)
的角动量相同（原点平移修正项
\(\boldsymbol a\times\boldsymbol P=0\)）。
碰前仅旋转着的两个质点有角动量，
\[
 L_{C,\rm before}
 =2m\left(\frac a2\right)^2\omega
 =\boxed{\frac12ma^2\omega}.
\]
碰撞冲量均为系统内力，若忽略外力矩则
\(\boxed{L_{C,\rm after}=L_{C,\rm before}}\)。
(3) 粘在碰撞端后该端质量是 \(2m\)，
距离新质心 \(a/3\)，另一端质量 \(m\)
距离新质心 \(2a/3\)；
故碰后转动惯量
\[
 I_C=2m(a/3)^2+m(2a/3)^2=\frac23ma^2.
\]
由 \(I_C\omega_f=L_C\) 得
\[
 \boxed{\omega_f=\frac{(ma^2\omega/2)}{(2ma^2/3)}
 =\frac34\,\omega}.
\]
系统总动能在非弹性碰撞中下降：
碰前 \(K_i=\frac14ma^2\omega^2\)，
碰后 \(K_f=\frac12I_C\omega_f^2
 =\frac3{16}ma^2\omega^2\)，符合耗散。
